\documentclass[aspectratio=169,12pt]{beamer}

\usepackage[T2A]{fontenc}
\usepackage[utf8]{inputenc}
\usepackage[english, main = russian]{babel}

\usepackage{graphicx}
\usepackage{amsmath, amssymb, amsfonts}
\usepackage{csquotes}
\usepackage{booktabs}

\usepackage{tikz}
\usepackage{pgfplots}
\pgfplotsset{compat=newest}
\pgfplotsset{
  nirsaxis/.style={
    width=0.85\textwidth,
    height=0.48\textwidth,
    tick label style={font=\scriptsize},
    label style={font=\small},
    xlabel style={yshift=-2pt},
    ylabel style={xshift=6pt},
  }
}

\usetheme{Madrid}
\setbeamertemplate{navigation symbols}{}

\title[Реакторные $\bar{\nu}_e$ в черенковском детекторе]{Регистрация реакторных антинейтрино\\с помощью черенковского детектора}
\subtitle{Geant4-моделирование prompt-сигнала IBD и сравнение допантов Cd/Gd}
\author{~Панфилов ~П.\,А.}
\institute{НИЯУ МИФИ}

\begin{document}

% -------------------------------------------------------------------------
\begin{frame}
  \titlepage

  \vspace{-0.4cm}
  \begin{center}
  \small
  \begin{tabular}{p{0.52\textwidth}p{0.42\textwidth}}
    Научный руководитель\\(старший преподаватель) &
    ~Мачулин ~И.\,Н.\\[0.35cm]
    Научный консультант\\(к.ф.-м.н.) &
    ~Долганов ~Г.\,Д.\\
  \end{tabular}
  \end{center}
\end{frame}

% -------------------------------------------------------------------------
\begin{frame}{Цель и задачи работы}
Реакторные антинейтрино $\bar{\nu}_e$ — интенсивный источник частиц с энергиями порядка нескольких MeV. Регистрация $\bar{\nu}_e$ в водных черенковских детекторах требует выделения как prompt-сигнала позитрона, так и delayed-сигнала захвата нейтрона.

\medskip
\textbf{Цель}: разработать Geant4-модель черенковского детектора для регистрации реакторных $\bar{\nu}_e$ и изучить отклик по свету для prompt- и delayed-компонент.

\medskip
В работе решались задачи:
\begin{itemize}
    \item разыгрывание спектра позитронов от реакции обратного бета-распада (IBD);
    \item учёт спектральных оптических свойств воды $n(\lambda)$, $L_{abs}(\lambda)$ и $QE(\lambda)$ ФЭУ;
    \item моделирование термализации и захвата тепловых нейтронов на допантах Cd и Gd;
    \item сравнение допантов по времени до захвата $t_{cap}$ и числу фотоэлектронов $N_{pe}$ от $\gamma$-каскада.
\end{itemize}
\end{frame}

% -------------------------------------------------------------------------
\begin{frame}{Физика сигнала: IBD prompt + delayed}
Классический канал регистрации реакторных антинейтрино:
\begin{equation}
    \bar{\nu}_e + p \rightarrow e^+ + n,\qquad T_{e^+} \approx E_{\bar{\nu}} - 1.806~\text{MeV}.
\end{equation}

\smallskip
Сигнал состоит из двух компонент:
\begin{itemize}
    \item \textit{prompt}: позитрон + аннигиляция $\gamma\gamma$ $\Rightarrow$ прямой черенковский свет;
    \item \textit{delayed}: термализация нейтрона ($\sim 1$–$10\,\mu$s) + захват на H, Cd или Gd $\Rightarrow$ $\gamma$-каскад $\Rightarrow$ комптоновские электроны $\Rightarrow$ вторичный черенковский свет.
\end{itemize}

\smallskip
\textbf{Энергетический бюджет «в свет»}:
\begin{center}
\small
\begin{tabular}{lcc}
\toprule
Компонента & Энергия & Множественность \\
\midrule
$e^+$ (prompt, средн.) & $T_{e^+} + 2\cdot 0{.}511 \approx 3{.}4$ MeV & --- \\
$\gamma$-каскад на $^{113}$Cd & $\approx 9$ MeV & $\bar M_\gamma \approx 2$–3 \\
$\gamma$-каскад на $^{157}$Gd & $\approx 8$ MeV & $\bar M_\gamma \approx 4$–5 \\
\bottomrule
\end{tabular}
\end{center}
\end{frame}

% -------------------------------------------------------------------------
\begin{frame}{Захват тепловых нейтронов на Cd и Gd}
После термализации на протонах воды нейтрон достигает тепловой энергии $E_{th}\sim k_BT\approx 0{.}025$ eV ($v_{th}\approx 2200$ m/s) и захватывается одним из ядер допанта.

\medskip
\textbf{Сечения радиационного захвата $(n,\gamma)$ на тепловых нейтронах:}
\begin{itemize}
    \item $^1$H: $\sigma\approx 0{.}33$ b ($E_\gamma = 2{.}22$ MeV);
    \item $^{113}$Cd: $\sigma\approx 2{.}0\cdot 10^4$ b ($\sum E_\gamma\approx 9$ MeV);
    \item $^{157}$Gd: $\sigma\approx 2{.}55\cdot 10^5$ b ($\sum E_\gamma\approx 8$ MeV).
\end{itemize}

\medskip
Среднее время до захвата:
\begin{equation}
    \tau_{cap} = \frac{1}{n_{dop}\,\sigma_{cap}\,v_{th}}.
\end{equation}

\smallskip
\textbf{Ожидание}: $\sigma_\text{Gd} \gg \sigma_\text{Cd}$ при той же массовой доле $\Rightarrow$ Gd должен давать значимо более короткое $\tau_{cap}$ и более узкое coincidence-окно для выделения IBD.
\end{frame}

% -------------------------------------------------------------------------
\begin{frame}{Черенковское излучение: что фиксирует детектор}
Условие возникновения и угол конуса:
\begin{equation}
    v>\frac{c}{n(\lambda)} \;\Leftrightarrow\; \beta n(\lambda)>1,\qquad \cos\theta_C(\lambda)=\frac{1}{\beta n(\lambda)}.
\end{equation}

\medskip
Спектр (формула Франка–Тамма):
\begin{equation}
    \frac{d^2N}{dx\,d\lambda}=2\pi\alpha\left(1-\frac{1}{\beta^2 n^2(\lambda)}\right)\frac{1}{\lambda^2}.
\end{equation}

\medskip
Для реалистичности модели критичны:
\begin{itemize}
    \item $n(\lambda)$ воды — задаёт порог и угол излучения;
    \item $L_{abs}(\lambda) = \lambda / (4\pi k(\lambda))$ — определяет дальность транспорта;
    \item $QE(\lambda)$ ФЭУ — вероятность преобразования $\gamma \to$ фотоэлектрон.
\end{itemize}
\end{frame}

% -------------------------------------------------------------------------
\begin{frame}{Geant4-модель: геометрия}
\begin{columns}[T]
\begin{column}{0.55\textwidth}
\textbf{Цилиндрическая компоновка}:
\begin{itemize}
    \item стальной бак $R = 630$ мм, $H = 1300$ мм, отражающие стенки ($R_{refl}=0{.}9$);
    \item внешний слой воды;
    \item внутренний PMMA-сосуд $R = 590$ мм, $H = 700$ мм;
    \item мишень — вода + допант (Cd или Gd, массовая доля $0{.}1\%$).
\end{itemize}
\end{column}
\begin{column}{0.45\textwidth}
\textbf{ФЭУ}:
\begin{itemize}
    \item всего \textbf{38} ФЭУ;
    \item по \textbf{19} на верхней и нижней крышках;
    \item схема концентрических колец $1+6+12$;
    \item радиус ФЭУ $100$ мм.
\end{itemize}

\medskip
\textbf{Переключение допанта}:\\ \texttt{/det/dopant Cd|Gd}.
\end{column}
\end{columns}

\smallskip
\textbf{Измеряемая величина}: суммарное число фотоэлектронов $N^{tot}_{pe}$ во всех ФЭУ за событие.
\end{frame}

% -------------------------------------------------------------------------
\begin{frame}{Geant4-модель: физика и генерация частиц}
\textbf{Подключённые конструкторы физики:}
\begin{itemize}
    \item \texttt{G4EmStandardPhysics} + \texttt{G4OpticalPhysics} (ЭМ и черенковский свет);
    \item \texttt{G4HadronPhysicsFTFP\_BERT\_HP}, \texttt{G4HadronElasticPhysicsHP},\\ \texttt{G4IonPhysics}, \texttt{G4DecayPhysics};
    \item опция \texttt{HP} (High Precision) использует данные G4NDL/ENDF-B — корректные тепловые сечения и резонансы $^{113}$Cd, $^{157}$Gd.
\end{itemize}

\medskip
\textbf{Два режима генерации первичных частиц:}
\begin{itemize}
    \item \texttt{/gen/mode ibd} — позитроны от IBD-спектра, изотропное направление, равномерное разыгрывание точки рождения по объёму мишени;
    \item \texttt{/gen/mode mono} \texttt{/gen/particle neutron} \texttt{/gen/monoEnergy 10 keV} — тепловые нейтроны, изотропно и равномерно по объёму, для исследования delayed-компоненты.
\end{itemize}
\end{frame}

% -------------------------------------------------------------------------
\begin{frame}{Результат 1: спектр кинетической энергии позитронов}
\begin{center}
\begin{tikzpicture}
\begin{axis}[
    nirsaxis,
    ybar,
    bar width=0.09,
    xlabel={$T_{e^+}$, MeV},
    ylabel={Counts},
    xmin=0, xmax=8,
    grid=both,
    minor grid style={gray!15},
    major grid style={gray!30},
]
\addplot[fill=blue!55, draw=blue!80]
table[x index=0, y index=1] {data/positron_Te_hist_trig3.dat};
\end{axis}
\end{tikzpicture}
\end{center}
\vspace{-0.4em}
\small Разыгрывание $E_{\bar{\nu}}$ по $\phi(E_{\bar{\nu}})\cdot\sigma_{\text{IBD}}(E_{\bar{\nu}})$, далее $T_{e^+}\approx E_{\bar{\nu}}-1.806$ MeV. $10^5$ событий, условие $N_{fired}\ge 3$: проходит $76.1\%$, $\langle T_{e^+}\rangle \approx 2{.}72$ MeV.
\end{frame}

% -------------------------------------------------------------------------
\begin{frame}{Результат 2: $\langle N_{pe}\rangle$ от энергии позитрона}
\begin{center}
\begin{tikzpicture}
\begin{axis}[
    nirsaxis,
    xlabel={$T_{e^+}$, MeV},
    ylabel={$\langle N^{tot}_{pe}\rangle$},
    xmin=0, xmax=8,
    grid=both,
    minor grid style={gray!15},
    major grid style={gray!30},
]
\addplot+[mark=*, mark size=1.2pt, line width=0.8pt, blue,
    error bars/.cd,
    y dir=both,
    y explicit
]
table[x index=0, y index=1, y error index=2] {data/npe_vs_Te_mean_trig3.dat};
\end{axis}
\end{tikzpicture}
\end{center}
\vspace{-0.4em}
\small $10^5$ событий IBD, условие $N_{fired}\ge 3$. Среднее по прошедшим триггер событиям: $\langle N^{tot}_{pe}\rangle \approx 23{.}4$.
\end{frame}

% -------------------------------------------------------------------------
\begin{frame}{Результат 3: время жизни нейтронов Cd vs Gd}
\begin{center}
\begin{tikzpicture}
\begin{axis}[
    nirsaxis,
    xlabel={$t_{cap}$, $\mu$s},
    ylabel={Counts},
    xmin=0, xmax=200,
    grid=both,
    minor grid style={gray!15},
    major grid style={gray!30},
    legend pos=north east,
    legend style={font=\scriptsize},
]
\addplot+[const plot, draw=blue, fill=blue!20, fill opacity=0.5, mark=none]
table[x index=0, y index=1] {data/lifetime_hist_Cd_trig3.dat};
\addlegendentry{Cd (0.1\%)}
\addplot+[const plot, draw=red, fill=red!20, fill opacity=0.5, mark=none]
table[x index=0, y index=1] {data/lifetime_hist_Gd_trig3.dat};
\addlegendentry{Gd (0.1\%)}
\end{axis}
\end{tikzpicture}
\end{center}
\vspace{-0.4em}
\small $T_n=10$ keV, $10^5$ событий, условие $N_{fired}\ge 3$. Средние времена: $\langle t_{cap}^{Cd}\rangle \approx 112{.}2\,\mu$s, $\langle t_{cap}^{Gd}\rangle \approx 31{.}6\,\mu$s.
\end{frame}

% -------------------------------------------------------------------------
\begin{frame}{Результат 4: $N_{pe}$ от захвата нейтрона Cd vs Gd}
\begin{columns}[T,onlytextwidth]
\begin{column}{0.5\textwidth}
\centering
\begin{tikzpicture}
\begin{axis}[
    nirsaxis,
    width=\textwidth,
    height=0.72\textwidth,
    title={\scriptsize Гистограмма $N^{tot}_{pe}$, log $Y$},
    title style={yshift=-1ex},
    xlabel={$N^{tot}_{pe}$},
    ylabel={Counts},
    xmin=0, xmax=300,
    ymode=log,
    ymin=1, ymax=80000,
    grid=both,
    minor grid style={gray!15},
    major grid style={gray!30},
    legend pos=north east,
    legend style={font=\tiny},
    legend cell align={left},
]
\addplot+[const plot mark left, draw=blue, mark=none, line width=0.9pt, fill=none]
table[x index=0, y index=1] {data/npe_hist_log_Cd_trig3.dat};
\addlegendentry{Cd (0.1\%)}
\addplot+[const plot mark left, draw=red, mark=none, line width=0.9pt, fill=none]
table[x index=0, y index=1] {data/npe_hist_log_Gd_trig3.dat};
\addlegendentry{Gd (0.1\%)}
\end{axis}
\end{tikzpicture}
\end{column}
\begin{column}{0.5\textwidth}
\centering
\begin{tikzpicture}
\begin{axis}[
    nirsaxis,
    width=\textwidth,
    height=0.72\textwidth,
    title={\scriptsize Кривая выживания $P(N^{tot}_{pe}\ge N^{tot}_{pe,thr})$},
    title style={yshift=-1ex},
    xlabel={$N^{tot}_{pe,thr}$},
    ylabel={$P(N^{tot}_{pe}\ge x)$},
    xmin=0, xmax=200,
    ymin=0, ymax=1,
    grid=both,
    minor grid style={gray!15},
    major grid style={gray!30},
    legend pos=north east,
    legend style={font=\tiny},
    legend cell align={left},
]
\addplot+[draw=blue, mark=none, line width=1.1pt]
table[x index=0, y index=1] {data/npe_survival_Cd_trig3.dat};
\addlegendentry{Cd (0.1\%)}
\addplot+[draw=red, mark=none, line width=1.1pt]
table[x index=0, y index=1] {data/npe_survival_Gd_trig3.dat};
\addlegendentry{Gd (0.1\%)}
\end{axis}
\end{tikzpicture}
\end{column}
\end{columns}
\vspace{0.3em}
{\scriptsize Условие $N_{fired}\ge 3$. $\langle N_{pe}^{Cd}\rangle\approx 40{.}6$, $\langle N_{pe}^{Gd}\rangle\approx 42{.}8$. Преимущество Gd сохраняется: при $N^{tot}_{pe,thr}=50$ --- $42\%$ vs $27\%$. Причина --- большая множественность $\gamma$-каскада ($\bar M_\gamma\approx 4$--5 у Gd vs $\sim 2$--3 у Cd).\par}
\end{frame}

% -------------------------------------------------------------------------
\begin{frame}{Условие регистрации: $N_{fired}\ge 3$}
\small
Тёмновые токи фотоумножителей ($r_{dark}\approx 1$~кГц/ФЭУ) дают суммарный фоновый поток $N_{PMT}\cdot r_{dark}=38$~кГц при одиночном пороге. При требовании $N_{fired}\ge k$ в окне $\tau_w$ частота случайных совпадений тёмновых токов:
\begin{equation*}
    R_{dark}^{(\ge k)}\approx\binom{N_{PMT}}{k}\,r_{dark}^{k}\,\tau_w^{k-1}
\end{equation*}
При $k=3$, $\tau_w=100$~нс: $R_{dark}^{(\ge 3)}\approx 0.08$~Гц. Все результаты выше получены с этим условием.

\medskip
\begin{center}
\begin{tabular}{lccc}
\toprule
Порог $N_{fired,thr}$ & IBD prompt & delayed Cd ($0.1\%$) & delayed Gd ($0.1\%$) \\
\midrule
$\ge 1$  & $90.0\%$ & $75.8\%$ & $87.5\%$ \\
$\ge 3$  & $76.1\%$ & $56.8\%$ & $74.0\%$ \\
$\ge 5$  & $61.6\%$ & $43.4\%$ & $61.6\%$ \\
$\ge 10$ & $28.5\%$ & $22.2\%$ & $34.3\%$ \\
\bottomrule
\end{tabular}
\end{center}

\medskip
{\small Кривые выживания $P(N_{fired}\ge N_{fired,thr})$ приведены в приложении.}
\end{frame}

% -------------------------------------------------------------------------
\begin{frame}{Выводы и дальнейшая работа}
\textbf{Итоги}:
\begin{itemize}\setlength{\itemsep}{1pt}
    \item реализована Geant4-модель черенковского детектора с prompt- и delayed-компонентами IBD;
    \item учтены спектральные свойства воды ($n(\lambda)$, $L_{abs}(\lambda)$) и $QE(\lambda)$ ФЭУ;
    \item prompt ($N_{fired}\ge 3$, $76.1\%$ событий): $\langle N^{tot}_{pe}\rangle\approx 23{.}4$, $\langle T_{e^+}\rangle\approx 2{.}72$ MeV;
    \item delayed ($N_{fired}\ge 3$): \textbf{Gd предпочтительнее Cd} по всем параметрам --- $t_{cap}$ в $3.5\times$ короче ($31.6$ vs $112.2\,\mu$s), $\langle N_{pe}\rangle$ выше ($42.8$ vs $40.6$), $\varepsilon_{trig}$ выше ($74\%$ vs $57\%$);
    \item тёмновые совпадения при $N_{fired}\ge 3$ и $\tau_w=100$~нс: $\approx 0.08$~Гц.
\end{itemize}

\smallskip
\textbf{Дальнейшая работа}:
\begin{itemize}\setlength{\itemsep}{1pt}
    \item учёт угловой корреляции позитрона с $\bar{\nu}_e$ (формализм Vogel–Beacom);
    \item объединение prompt+delayed в IBD-цепочку с триггером совпадений;
    \item оценка ожидаемой скорости счёта на Калинской АЭС.
\end{itemize}
\end{frame}

% =========================================================================
\appendix
% =========================================================================

\begin{frame}{Приложение: кривые выживания по $N_{fired}$}
\begin{center}
\begin{tikzpicture}
\begin{axis}[
    nirsaxis,
    width=0.85\textwidth,
    height=0.46\textwidth,
    xlabel={$N_{fired,thr}$ (порог по числу сработавших ФЭУ)},
    ylabel={$P(N_{fired}\ge N_{fired,thr})$},
    xmin=0, xmax=20,
    ymin=0, ymax=1,
    grid=both,
    minor grid style={gray!15},
    major grid style={gray!30},
    legend pos=north east,
    legend style={font=\scriptsize},
    legend cell align={left},
]
\addplot+[draw=blue!80!black, mark=none, line width=1.1pt]
table[x index=0, y index=1] {data/nfired_survival_ibd.dat};
\addlegendentry{IBD prompt}
\addplot+[draw=red, mark=none, line width=1.1pt]
table[x index=0, y index=1] {data/nfired_survival_Gd.dat};
\addlegendentry{delayed Gd (0.1\%)}
\addplot+[draw=teal, mark=none, line width=1.1pt]
table[x index=0, y index=1] {data/nfired_survival_Cd.dat};
\addlegendentry{delayed Cd (0.1\%)}
\addplot+[draw=gray, dashed, mark=none, line width=0.8pt]
coordinates {(3,0) (3,1)};
\end{axis}
\end{tikzpicture}
\end{center}
\vspace{-0.4em}
\small Пунктир --- выбранный рабочий порог $N_{fired}=3$.
\end{frame}

\end{document}
